% =============================================================================
%  3-Conceptos.tex — Masking, hiperfoco y la tríada burnout/meltdown/shutdown
% =============================================================================
%  NOTA: citas en traducción propia; original en comentario (% ORIG:).
%  EDITAR: falta número de página exacto en citas tomadas del texto completo
%          de PMC (Hull et al., Raymaker et al., Ashinoff y Abu-Akel) y del
%          manuscrito aceptado de Wood (paginación distinta a la publicada).
%          Phung et al. y Hours et al. son artículos con número de artículo,
%          sin paginación.
% =============================================================================

\section{Desarrollo teórico: conceptos clave}

Los conceptos de esta sección se presentan, en la medida de lo posible, desde las descripciones de las propias personas autistas recogidas en la investigación, y no desde un listado de síntomas.

\subsection{Masking (enmascaramiento)}

\citet[p.~2519]{hull2017} señalan que \enquote{se plantea como hipótesis que el camuflaje de las características autistas en situaciones sociales es una estrategia común de afrontamiento social en personas adultas con condiciones del espectro autista}, y que \enquote{el camuflaje está motivado por el deseo de encajar con los demás y de establecer vínculos}.
% ORIG: Camouflaging of autistic characteristics in social situations is
%       hypothesised as a common social coping strategy for adults with autism
%       spectrum conditions (ASC).
% ORIG: Camouflaging is motivated by the desire to fit in with others and to
%       make connections.
Entre las técnicas reportadas, una es central para esta sesión: \enquote{forzar y mantener un contacto visual apropiado, o intentar mirar lo más cerca posible de los ojos del otro, fue una técnica compensatoria comúnmente reportada} \citep{hull2017}.
% ORIG: Forcing and maintaining appropriate eye contact, or attempting to look as
%       close to another's eyes as possible, was a common compensatory technique
%       reported.

El costo es alto: \enquote{con mucho, la consecuencia más constante del camuflaje descrita por los participantes fue el agotamiento. El camuflaje fue descrito con frecuencia como algo mental, física y emocionalmente agotador, que requiere una concentración intensa, autocontrol y manejo de la incomodidad} \citep{hull2017}. Un testimonio lo expresa así: \enquote{es mentalmente agotador tener que ser constantemente otra cosa, literalmente no poder ser nunca yo mismo} \citep{hull2017}.
% ORIG: By far the most consistent consequence of camouflaging described by
%       respondents was exhaustion. Camouflaging was frequently described as
%       being mentally, physically, and emotionally draining; requiring
%       intensive concentration, self-control, and management of discomfort.
% ORIG: It's mentally exhausting constantly having to be something else,
%       literally never being able to be myself

\subsection{Hiperfoco}

\citet{ashinoff2021} proponen una definición operativa con cuatro rasgos: \enquote{(1) el hiperfoco es inducido por la implicación en una tarea; (2) el hiperfoco se caracteriza por un estado intenso de atención sostenida o selectiva; (3) durante un estado de hiperfoco hay una percepción disminuida de los estímulos no relevantes para la tarea, y (4) durante un estado de hiperfoco, el desempeño en la tarea mejora}.
% ORIG: (1) hyperfocus is induced by task engagement; (2) hyperfocus is
%       characterized by an intense state of sustained or selective attention;
%       (3) During a hyperfocus state, there is a diminished perception of
%       non-task relevant stimuli; and (4) During a hyperfocus state, task
%       performance improves.

En el contexto escolar, \citet[p.~34]{wood2021} lo aborda desde los intereses intensos: \enquote{tener intereses intensos o \enquote{especiales}, y una tendencia a concentrarse en profundidad con exclusión de otros estímulos, se asocia con la cognición autista, a veces enmarcada como \enquote{monotropismo}. A pesar de algunos inconvenientes y de asociaciones negativas con la repetición no deseada, esta disposición se vincula con una serie de beneficios educativos y de largo plazo para los niños autistas}. De ahí, concluye la autora, \enquote{la importancia de percibir esta disposición como una fortaleza, en lugar de un déficit que deba remediarse de algún modo} \citep{wood2021}.
% ORIG: Having intense or "special" interests and a tendency to focus in depth
%       to the exclusion of other inputs, is associated with autistic
%       cognition, sometimes framed as "monotropism". Despite some drawbacks and
%       negative associations with unwanted repetition, this disposition is
%       linked to a range of educational and longer-term benefits for autistic
%       children.
% ORIG: the importance of perceiving this disposition as a strength
%       (Winter-Messiers 2007) rather than a deficit which must somehow be
%       remedied

El hiperfoco también se describe en el TDAH, condición que con frecuencia acompaña al autismo. \citet{hours2022} indican que, \enquote{según la literatura científica, entre el 50 y el 70\,\% de las personas con trastorno del espectro autista (TEA) presentan también trastorno por déficit de atención e hiperactividad (TDAH) comórbido}, pero aclaran que la nosografía actual \enquote{deja claro que se trata de entidades diferentes, que afectan a los niños y a sus historias de desarrollo de maneras claramente distintas}. Por eso suelen estudiarse a la par, aunque no significan lo mismo. Sobre el TDAH, \citet{ashinoff2021} señalan que, \enquote{aunque el hiperfoco parece contradecir la asociación del TDAH con la inatención y la impulsividad, a menudo se reporta como un estado positivo en personas con TDAH, porque de hecho se implican en tareas durante periodos más largos de lo habitual}.
% ORIG: According to the scientific literature, 50 to 70% of individuals with
%       autism spectrum disorder (ASD) also present with comorbid attention
%       deficit hyperactivity disorder (ADHD).
% ORIG: [The semiology of ASD and ADHD presented in current nosography ...]
%       makes it clear that they are different entities, affecting children
%       and their developmental histories in ways that are clearly distinct.
% ORIG: although hyperfocus is seemingly antithetical to the association of ADHD
%       with inattention and impulsivity, it is often reported as a positive
%       state in individuals with ADHD because they actually engage in tasks
%       for longer periods of time than is typical

\subsection{La tríada del estrés y la regulación}

\citet{phung2021} agrupan cuatro experiencias: \enquote{el burnout, la inercia, el meltdown y el shutdown (BIMS) se han identificado como partes importantes de la vida de algunas personas autistas}. Esta guía retoma tres de ellas.
% ORIG: Burnout, inertia, meltdown, and shutdown (BIMS) have been identified as
%       important parts of some autistic people's lives.

\paragraph{Burnout.} \citet{raymaker2020} lo definen así: \enquote{el burnout autista es un síndrome conceptualizado como resultado del estrés vital crónico y de un desajuste entre expectativas y capacidades sin los apoyos adecuados. Se caracteriza por un agotamiento generalizado y de largo plazo (típicamente de tres meses o más), pérdida de funciones y menor tolerancia a los estímulos}. Entre sus causas, el estudio identifica \enquote{los estresores de la vida (subtemas: enmascaramiento, expectativas, manejo de la discapacidad y cambios vitales) que conducen a una carga acumulada} \citep{raymaker2020}.
% ORIG: Autistic burnout is a syndrome conceptualized as resulting from chronic
%       life stress and a mismatch of expectations and abilities without
%       adequate supports. It is characterized by pervasive, long-term
%       (typically 3+ months) exhaustion, loss of function, and reduced
%       tolerance to stimulus.
% ORIG: life stressors (subthemes: masking, expectations, disability
%       management, and life-changes) that lead to (2) cumulative load

\paragraph{Meltdown.} A partir de investigaciones previas con informantes autistas, \citet{phung2021} lo describen como \enquote{un fenómeno con expresiones variadas en el que los informantes autistas se sienten completamente abrumados, acompañado de una falta de control y de estrés acumulado}, que provoca \enquote{respuestas de ansiedad hacia el exterior y descarga de energía}. Entre sus detonantes mencionan \enquote{demandas sociales, frustración, vergüenza, dificultades de comunicación, detonantes emocionales y estímulos sensoriales aversivos abrumadores}. Una de las personas jóvenes entrevistadas lo describe así: \enquote{eres un pasajero en un viaje de destrucción… y es como ir chocando con un montón de cosas} \citep{phung2021}.
% ORIG: A phenomenon with varying expressions by which autistic informants feel
%       entirely overwhelmed accompanied by a lack of control and cumulative
%       stress
% ORIG: Meltdowns elicit responses of outward anxiety and energy outpour
% ORIG: social demands, frustration, embarrassment, challenges with
%       communication, emotional triggers, and overwhelming aversive sensory
%       stimuli
% ORIG: You're a passenger on a ride of destruction… and it's like hitting a
%       bunch of stuff.

\paragraph{Shutdown.} \enquote{Aunque similares a los meltdowns, los shutdowns se presentan como experiencias más internas, en las que la persona se retira de su entorno y que van acompañadas de dolor emocional} \citep{phung2021}. Su intensidad varía: \enquote{va de leve (p.\,ej., poder caminar y hablar) a severo (p.\,ej., sentirse desconectado de las propias extremidades y adoptar una posición fetal)}. Los autores retoman a Belek, cuya descripción \enquote{enfatiza detonantes como la sobrecarga sensorial y experiencias internas de miedo y parálisis}. En palabras de una persona joven: \enquote{el shutdown es, como, normalmente por estrés: me congelo por completo, no puedo hablar bien, tartamudeo, me congelo literalmente} \citep{phung2021}.
% ORIG: Although similar to meltdowns, shutdowns present as more internal
%       experiences, where the individual withdraws from their surroundings and
%       is accompanied by emotional pain
% ORIG: The degree to which one can function during a shutdown ranges from mild
%       (e.g., being able to walk around and talk) to severe (e.g., feeling
%       detached from your limbs and going into a fetal position)
% ORIG: Belek's description of shutdown emphasizes triggers like sensory
%       overload and internal experiences of fear and paralysis
% ORIG: Shutdown is just, like, usually 'cause of stress, I'll just completely
%       freeze up, I can't really talk, I'll stutter, I'll like literally
%       freeze up.
